\subsection{Monte-Carlo checks}\label{sec:mc}

Table~\ref{tab:mc} uses a synthetic MNAR model with known propensities: 30 users,
40 items, low-rank preferences, and exposures driven by user activity, item
popularity and the outcome itself (median propensity 0.06, a third of the
pairs below 0.05). We draw 20{,}000 exposure matrices for three hops and 4{,}000
for five hops.

\begin{table}[t]\centering\small
\caption{Monte-Carlo check of Theorems~\ref{thm:unbiased}, \ref{thm:khop}, \ref{thm:lower}
and~\ref{thm:fair} and Corollary~\ref{cor:cand}. Relative $|$bias$|$ is the mean absolute
bias divided by the mean target, over all entries and conditionally on unexposed
candidates (against the conditional target of Corollary~\ref{cor:cand}). Kendall
$\tau$ is the mean within-user rank agreement of one realisation with the
full-exposure target on unexposed candidates. $\eta$ is the exposure elasticity
(separate model, 3{,}000 replications).}\label{tab:mc}
\begin{adjustbox}{max width=\textwidth}
\begin{tabular}{lcccc}
\toprule
Estimator & rel.\,|bias| & frac.\ $|z|>3$ & rel.\,RMSE & exposure elasticity $\eta$\\
\midrule
Obs & 0.992 & 1.0000 & 1.02 & +1.777\\
IPS & 2.403 & 0.4042 & 21.01 & +0.007\\
IPS+WC & 0.028 & 0.0025 & 4.93 & +0.007\\
DR & 1.903 & 0.6292 & 20.94 & -0.187\\
DRUP & 0.020 & 0.0008 & 3.97 & +0.005\\
\bottomrule
\end{tabular}

\end{adjustbox}
\end{table}

Over all entries, the uncorrected IPS and DR operators are biased by 2.3 and 1.8
times the mean target at three hops, and by 15 and 8 times at five hops. With
the correction the residual bias (0.02--0.06) is at the level of the
Monte-Carlo error. On unexposed candidates the picture is the one predicted by
Corollary~\ref{cor:cand}: the uncorrected IPS operator is unbiased at three
hops (0.029, identical to IPS with the correction) but not at five hops (4.7),
and the uncorrected DR operator is biased at both depths (1.1 and 5.1). The
exposure elasticity is $+1.78$ for the logged graph, $-0.19$ for DR adjacency and
$+0.005$ for DRUP. The $K$-hop engine reproduces the closed-form three-hop
operator to $7\times10^{-15}$ and brute-force enumeration of all five-hop walks to
$10^{-13}$.

Unbiasedness does not buy ranking accuracy in this model. The rank agreement
of a single realisation with the target is higher for the uncorrected DR
operator (Kendall $\tau=0.20$) than for DRUP (0.12), and at five hops 0.11 against
0.09. The reason is visible in Theorem~\ref{thm:lower}: on unexposed candidates
the bias is $C_{ui}\hat Y_{ui}(R_u^{(i)}+K_i)$, which adds a positive multiple of the
imputation to the score, and here the imputation is informative. The corrected
operator removes this term and, with it, some useful signal. For IPS edges,
whose bias does not involve the imputation, the correction leaves the rank
agreement unchanged (0.061 against 0.059 at five hops).

\begin{table}[t]\centering\small
\caption{Stress tests outside Assumption~\ref{as:unconf} and with misspecified
propensities (log-normal noise with $\sigma=0.5$). Relative $|$bias$|$ on unexposed
candidates.}\label{tab:stress}
\begin{adjustbox}{max width=\textwidth}
\begin{tabular}{lcccc}
\toprule
Scenario ($K=3$, rel.\ $|$bias$|$ on unexposed candidates) & IPS & IPS + WC & DR & \textbf{DRUP}\\
\midrule
independent exposures (Assumption~\ref{as:unconf}) & 0.029 & 0.029 & 1.137 & 0.019\\
fixed-size slates per user (Proposition~\ref{prop:dep}) & 0.219 & 0.219 & 1.119 & 0.024\\
misspecified propensities, exact imputation & 0.568 & 0.568 & 0.000 & 0.000\\
misspecified propensities, noisy imputation & 0.568 & 0.568 & 1.916 & 0.158\\
independent exposures (Assumption~\ref{as:unconf}), $K=5$ & 4.667 & 0.059 & 5.053 & 0.040\\
fixed-size slates per user (Proposition~\ref{prop:dep}), $K=5$ & 2.392 & 0.416 & 4.633 & 0.032\\
misspecified propensities, noisy imputation, $K=5$ & 19.266 & 1.162 & 16.027 & 0.204\\
\bottomrule
\end{tabular}

\end{adjustbox}
\end{table}

Table~\ref{tab:stress} leaves the assumptions of Theorem~\ref{thm:unbiased}. With
fixed-size slates (every user sees exactly $k_u$ items, drawn without
replacement, so exposures within a user are negatively dependent with
$\rho\approx0.29$), IPS picks up a first-order bias (0.22 at three hops), whereas
DRUP stays at 0.024, as Proposition~\ref{prop:dep} predicts for DR edges. With
misspecified propensities and an exact imputation, both DR operators remain
exactly unbiased, the double robustness of Theorem~\ref{thm:unbiased}; with a
noisy imputation DRUP's bias is 0.16, against 0.57 for IPS and 1.9 for the
uncorrected DR operator. Finally, the variance and concentration bounds of
Theorem~\ref{thm:var} were never violated, but they are loose
(Table~\ref{tab:bounds}): by a factor of $5\cdot10^2$--$8\cdot10^3$ for the variance and
$7\cdot10^1$--$2\cdot10^3$ for the deviation at level 0.05. They describe how the error
scales with $\tau$ and the imputation error, not its size.

\begin{table}[t]\centering\small
\caption{Tightness of Theorem~\ref{thm:var} in a 12-user, 15-item model (20{,}000
replications, all propensities $\ge\tau$), over all $(u,i)$.}\label{tab:bounds}
\begin{tabular}{lcc}
\toprule
 & $\tau=0.3$ & $\tau=0.1$\\
\midrule
variance bound (b) / Monte-Carlo variance, median & 1582 & 8422\\
variance bound (b) / Monte-Carlo variance, minimum & 471 & 2126\\
McDiarmid half-width (d) / empirical 95\% quantile, median & 131 & 1831\\
McDiarmid half-width (d) / empirical 95\% quantile, minimum & 73 & 685\\
\bottomrule
\end{tabular}

\end{table}
